% ============================================================
% Introduction Générale
% ============================================================
\chapter*{Introduction Générale}
\addcontentsline{toc}{chapter}{Introduction Générale}

Le secteur des télécommunications repose sur une relation client de masse. Un opérateur traite chaque jour un volume considérable de sollicitations, dont la très grande majorité porte sur des opérations courantes : consulter un solde, régler une facture, recharger une ligne, activer une option, signaler une connexion dégradée ou débloquer une carte SIM. Ces demandes suivent un cheminement stable, s'appuient sur des informations déjà détenues par l'opérateur et aboutissent à un nombre restreint d'actions bien identifiées. Elles mobilisent pourtant une part importante du temps des conseillers, tandis que les canaux automatisés existants, construits autour d'arborescences de menus, imposent au client de traduire son besoin dans le vocabulaire du système et ne transportent pratiquement aucun contexte d'une étape à la suivante. Les progrès récents des modèles de langue et des technologies vocales rendent aujourd'hui envisageable une prise en charge conversationnelle de ces demandes, à condition de résoudre deux difficultés simultanément : tenir le rythme d'un échange parlé, et garantir qu'un système capable d'agir sur un compte client ne le fasse jamais sans autorisation vérifiable.

Le présent projet, réalisé au sein de la société Amsys Consulting, porte sur la conception et la réalisation d'une plateforme conversationnelle temps réel destinée à l'automatisation de la relation client dans ce secteur. La plateforme prend en charge les demandes fréquentes de bout en bout, en voix comme en texte et dans trois langues, en s'appuyant sur un ensemble d'agents spécialisés. Elle ancre ses réponses factuelles dans la documentation propre à l'opérateur plutôt que dans la mémoire du modèle, soumet toute action engageant un compte à des règles métier déterministes produisant un verdict motivé, exécute les opérations approuvées de manière à ce qu'elles ne puissent être réalisées deux fois, consigne chaque acte conséquent dans un registre dont l'altération est détectable, et transmet la conversation à un conseiller humain, accompagnée du dossier complet, lorsque l'autonomie n'est pas appropriée. Deux interfaces web complètent l'ensemble : un portail destiné au client et une console de supervision et d'administration destinée aux équipes internes.

Ce rapport est organisé en quatre chapitres. Le premier situe le projet dans son cadre, présente l'organisme d'accueil, expose la problématique et les objectifs retenus, examine les solutions existantes et décrit la chronologie ainsi que le processus de développement appliqué. Le deuxième traduit ces objectifs en spécifications : acteurs, besoins fonctionnels, besoins non fonctionnels et cas d'utilisation. Le troisième présente la conception du système, sa vue statique et sa vue dynamique, puis détaille la conception de la couche agentique et celle du contrôle d'accès. Le quatrième expose la réalisation : technologies employées, chaîne de traitement temps réel étudiée étape par étape, gestion de la concurrence, architecture de déploiement, interfaces produites, campagne de vérification, résultats obtenus et difficultés rencontrées. Une conclusion générale dresse le bilan du travail accompli et présente les perspectives d'évolution de la plateforme.
